\unitchapter{Paper 1 · Unit 6}{Logical Reasoning}{p1-06}

\textbf{Expected questions: \textasciitilde5 (10 marks) · Target: 5/5}

\begin{sylbox}

\begin{itemize}
\tightlist
\item[$\square$]
  Structure of arguments: argument forms, categorical propositions, mood \& figure, formal \& informal fallacies, uses of language, connotations \& denotations, classical square of opposition
\item[$\square$]
  Evaluating \& distinguishing deductive and inductive reasoning
\item[$\square$]
  Analogies
\item[$\square$]
  Venn diagram: simple \& multiple use for establishing validity of arguments
\item[$\square$]
  Indian Logic: means of knowledge (Pramanas), structure \& kinds of Anumana (inference), Hetvabhasas (fallacies of inference)
\end{itemize}

\end{sylbox}

\needspace{170pt}

\hypertarget{p1-06--1-argument-basics}{%
\section{1. Argument Basics}\label{p1-06--1-argument-basics}}

An \textbf{argument} = premises + conclusion. Premise indicators: \emph{since, because, for}. Conclusion indicators: \emph{therefore, hence, thus, so}.

\begingroup\small

\begin{longtable}[]{@{}
  >{\raggedright\arraybackslash}p{(\columnwidth - 4\tabcolsep) * \real{0.0947}}
  >{\raggedright\arraybackslash}p{(\columnwidth - 4\tabcolsep) * \real{0.3791}}
  >{\raggedright\arraybackslash}p{(\columnwidth - 4\tabcolsep) * \real{0.3587}}@{}}
\toprule\noalign{}
\begin{minipage}[b]{\linewidth}\raggedright
\end{minipage} & \begin{minipage}[b]{\linewidth}\raggedright
\textbf{Deductive}
\end{minipage} & \begin{minipage}[b]{\linewidth}\raggedright
\textbf{Inductive}
\end{minipage} \\
\midrule\noalign{}
\endhead
\bottomrule\noalign{}
\endlastfoot
Direction & General → particular (not strictly) & Particular → general \\
Conclusion & Necessarily follows & Probably follows \\
Evaluation & \textbf{Valid / invalid}; sound = valid + true premises & \textbf{Strong / weak}; cogent = strong + true premises \\
Example & All men are mortal; Socrates is a man; ∴ mortal & Every crow seen is black ∴ all crows are black \\
\end{longtable}

\endgroup

Validity is about \textbf{form}, not truth. An argument with false premises can be valid.

\needspace{135pt}

\hypertarget{p1-06--2-categorical-propositions-a-e-i-o}{%
\section{2. Categorical Propositions (A E I O)}\label{p1-06--2-categorical-propositions-a-e-i-o}}

\begingroup\small

\begin{longtable}[]{@{}
  >{\raggedright\arraybackslash}p{(\columnwidth - 8\tabcolsep) * \real{0.0588}}
  >{\raggedright\arraybackslash}p{(\columnwidth - 8\tabcolsep) * \real{0.1572}}
  >{\raggedright\arraybackslash}p{(\columnwidth - 8\tabcolsep) * \real{0.0959}}
  >{\raggedright\arraybackslash}p{(\columnwidth - 8\tabcolsep) * \real{0.1088}}
  >{\raggedright\arraybackslash}p{(\columnwidth - 8\tabcolsep) * \real{0.1148}}@{}}
\toprule\noalign{}
\begin{minipage}[b]{\linewidth}\raggedright
\textbf{Code}
\end{minipage} & \begin{minipage}[b]{\linewidth}\raggedright
\textbf{Form}
\end{minipage} & \begin{minipage}[b]{\linewidth}\raggedright
\textbf{Quantity}
\end{minipage} & \begin{minipage}[b]{\linewidth}\raggedright
\textbf{Quality}
\end{minipage} & \begin{minipage}[b]{\linewidth}\raggedright
\textbf{Distributes}
\end{minipage} \\
\midrule\noalign{}
\endhead
\bottomrule\noalign{}
\endlastfoot
\textbf{A} & All S are P & Universal & Affirmative & S only \\
\textbf{E} & No S is P & Universal & Negative & S and P \\
\textbf{I} & Some S are P & Particular & Affirmative & None \\
\textbf{O} & Some S are not P & Particular & Negative & P only \\
\end{longtable}

\endgroup

Memory: \textbf{A}ff\textbf{I}rmo (A, I affirmative), n\textbf{E}g\textbf{O} (E, O negative). Distribution: "\textbf{ASEBINOP}" --- A: Subject, E: Both, I: Nothing, O: Predicate.

\needspace{186pt}

\hypertarget{p1-06--square-of-opposition}{%
\subsection{Square of Opposition}\label{p1-06--square-of-opposition}}

\begin{asciiblock}{131pt}
\begin{Verbatim}[fontsize=\fontsize{8.8}{10.7}\selectfont]
        A ─────────── contraries ─────────── E
        │  ╲                               ╱  │
        │    ╲                           ╱    │
 sub-   │      ╲   contradictories     ╱      │  sub-
 alter- │        ╲                   ╱        │  alter-
 nation │          ╳               ╳          │  nation
        │        ╱                   ╲        │
        │      ╱                       ╲      │
        │    ╱                           ╲    │
        I ──────────── sub-contraries ──────── O
\end{Verbatim}
\end{asciiblock}

\begingroup\small

\begin{longtable}[]{@{}
  >{\raggedright\arraybackslash}p{(\columnwidth - 4\tabcolsep) * \real{0.1291}}
  >{\raggedright\arraybackslash}p{(\columnwidth - 4\tabcolsep) * \real{0.1053}}
  >{\raggedright\arraybackslash}p{(\columnwidth - 4\tabcolsep) * \real{0.5383}}@{}}
\toprule\noalign{}
\begin{minipage}[b]{\linewidth}\raggedright
\textbf{Relation}
\end{minipage} & \begin{minipage}[b]{\linewidth}\raggedright
\textbf{Pair}
\end{minipage} & \begin{minipage}[b]{\linewidth}\raggedright
\textbf{Rule}
\end{minipage} \\
\midrule\noalign{}
\endhead
\bottomrule\noalign{}
\endlastfoot
Contradictory & A--O, E--I & Exactly one is true (opposite truth values always) \\
Contrary & A--E & Both cannot be true; both can be false \\
Sub-contrary & I--O & Both cannot be false; both can be true \\
Sub-alternation & A→I, E→O & If universal true ⇒ particular true; if particular false ⇒ universal false \\
\end{longtable}

\endgroup

\textbf{Inference table} (if given proposition is TRUE):

\begingroup\small

\begin{longtable}[]{@{}
  >{\raggedright\arraybackslash}p{(\columnwidth - 8\tabcolsep) * \real{0.1061}}
  >{\raggedright\arraybackslash}p{(\columnwidth - 8\tabcolsep) * \real{0.0311}}
  >{\raggedright\arraybackslash}p{(\columnwidth - 8\tabcolsep) * \real{0.0311}}
  >{\raggedright\arraybackslash}p{(\columnwidth - 8\tabcolsep) * \real{0.0311}}
  >{\raggedright\arraybackslash}p{(\columnwidth - 8\tabcolsep) * \real{0.0311}}@{}}
\toprule\noalign{}
\begin{minipage}[b]{\linewidth}\raggedright
\textbf{Given true}
\end{minipage} & \begin{minipage}[b]{\linewidth}\raggedright
\textbf{A}
\end{minipage} & \begin{minipage}[b]{\linewidth}\raggedright
\textbf{E}
\end{minipage} & \begin{minipage}[b]{\linewidth}\raggedright
\textbf{I}
\end{minipage} & \begin{minipage}[b]{\linewidth}\raggedright
\textbf{O}
\end{minipage} \\
\midrule\noalign{}
\endhead
\bottomrule\noalign{}
\endlastfoot
A & T & F & T & F \\
E & F & T & F & T \\
I & ? & F & T & ? \\
O & F & ? & ? & T \\
\end{longtable}

\endgroup

If given FALSE: A false → O true, E?, I?; E false → I true; I false → E true, A false, O true; O false → A true, E false, I true.

\needspace{168pt}

\hypertarget{p1-06--3-categorical-syllogism--mood--figure}{%
\section{3. Categorical Syllogism --- Mood \& Figure}\label{p1-06--3-categorical-syllogism--mood--figure}}

Three terms: Major (P -- predicate of conclusion), Minor (S -- subject of conclusion), Middle (M -- in both premises, not conclusion).

\begin{asciiblock}{78pt}
\begin{Verbatim}[fontsize=\fontsize{8.8}{10.7}\selectfont]
Figure 1      Figure 2      Figure 3      Figure 4
M – P         P – M         M – P         P – M
S – M         S – M         M – S         M – S
─────         ─────         ─────         ─────
S – P         S – P         S – P         S – P
\end{Verbatim}
\end{asciiblock}

\textbf{Mood} = the letters of the three propositions, e.g. AAA-1 (Barbara), EAE-1 (Celarent), AII-1 (Darii), EIO-1 (Ferio).

\needspace{100pt}

\hypertarget{p1-06--rules-of-validity}{%
\subsection{Rules of validity}\label{p1-06--rules-of-validity}}

\begin{enumerate}
\def\labelenumi{\arabic{enumi}.}
\tightlist
\item
  Middle term must be distributed at least once (else \textbf{fallacy of undistributed middle}).
\item
  A term distributed in conclusion must be distributed in premise (else \textbf{illicit major/minor}).
\item
  Two negative premises → no conclusion (\textbf{exclusive premises}).
\item
  A negative premise → negative conclusion; and vice-versa.
\item
  Two particular premises → no conclusion.
\item
  Exactly three terms (else \textbf{four-term fallacy}).
\end{enumerate}

\needspace{141pt}

\hypertarget{p1-06--venn-diagram-method-for-syllogism-questions}{%
\subsection{Venn Diagram Method (for syllogism questions)}\label{p1-06--venn-diagram-method-for-syllogism-questions}}

\begin{asciiblock}{86pt}
\begin{Verbatim}[fontsize=\fontsize{8.5}{10.3}\selectfont]
 Premise: All A are B       Premise: Some B are C      Possible conclusions:
     ┌───────────┐                                    "Some A are C"  → NOT definite
     │  B  ┌───┐ │   ┌──────┐                         "Some C are B"  → TRUE (conversion of I)
     │     │ A │ │◄─►│  C   │
     │     └───┘ │   └──────┘
     └───────────┘
\end{Verbatim}
\end{asciiblock}

Technique: draw the \textbf{minimum overlap} diagram; a conclusion follows only if it is true in \textbf{every} possible diagram.

\needspace{250pt}

\hypertarget{p1-06--4-fallacies}{%
\section{4. Fallacies}\label{p1-06--4-fallacies}}

\needspace{210pt}

\hypertarget{p1-06--formal-fallacies}{%
\subsection{Formal fallacies}\label{p1-06--formal-fallacies}}

Undistributed middle, illicit major, illicit minor, affirming the consequent (p→q, q ∴ p), denying the antecedent (p→q, ¬p ∴ ¬q), four terms.

\needspace{135pt}

\hypertarget{p1-06--informal-fallacies}{%
\subsection{Informal fallacies}\label{p1-06--informal-fallacies}}

\begingroup\small

\begin{longtable}[]{@{}
  >{\raggedright\arraybackslash}p{(\columnwidth - 2\tabcolsep) * \real{0.2934}}
  >{\raggedright\arraybackslash}p{(\columnwidth - 2\tabcolsep) * \real{0.3104}}@{}}
\toprule\noalign{}
\begin{minipage}[b]{\linewidth}\raggedright
\textbf{Fallacy}
\end{minipage} & \begin{minipage}[b]{\linewidth}\raggedright
\textbf{Meaning}
\end{minipage} \\
\midrule\noalign{}
\endhead
\bottomrule\noalign{}
\endlastfoot
Ad hominem & Attacking the person, not the argument \\
Ad populum (bandwagon) & Many people believe it ⇒ true \\
Ad verecundiam & Appeal to inappropriate authority \\
Ad ignorantiam & Not proved false ⇒ true \\
Ad misericordiam & Appeal to pity \\
Ad baculum & Appeal to force \\
Straw man & Distorting opponent\textquotesingle s argument \\
Red herring & Diverting to an irrelevant issue \\
Begging the question (petitio principii) & Conclusion assumed in premise (circular) \\
Hasty generalisation & Too small a sample \\
False cause (post hoc) & After this, therefore because of this \\
Slippery slope & One step leads to disaster chain \\
False dilemma & Only two options presented \\
Equivocation & Word used in two senses \\
Composition / Division & Parts ↔ whole wrongly transferred \\
\end{longtable}

\endgroup

\needspace{100pt}

\hypertarget{p1-06--5-uses-of-language-connotation--denotation}{%
\section{5. Uses of Language; Connotation \& Denotation}\label{p1-06--5-uses-of-language-connotation--denotation}}

\begin{itemize}
\tightlist
\item
  Functions of language (Copi): \textbf{Informative, Expressive, Directive}, (also ceremonial, performative).
\item
  \textbf{Denotation (extension)} = the set of objects a term refers to. \textbf{Connotation (intension)} = properties/attributes.
\item
  As connotation increases, denotation decreases (inverse relation). e.g. "animal" → "mammal" → "dog".
\end{itemize}

\needspace{135pt}

\hypertarget{p1-06--6-analogies}{%
\section{6. Analogies}\label{p1-06--6-analogies}}

A : B :: C : ? --- identify relation: part-whole, cause-effect, tool-worker, synonym, antonym, product-raw material, degree.

\begin{itemize}
\tightlist
\item
  Pen : Writer :: Scalpel : \textbf{Surgeon} · Cow : Calf :: Horse : \textbf{Foal} · Ornithology : Birds :: Entomology : \textbf{Insects}
\end{itemize}

\needdiagram{figures/p1-06-00}{95pt}

\hypertarget{p1-06--7-indian-logic-nyaya--very-high-weightage}{%
\section{7. Indian Logic (Nyaya) --- Very High Weightage}\label{p1-06--7-indian-logic-nyaya--very-high-weightage}}

\needspace{100pt}

\hypertarget{p1-06--pramanas-means-of-valid-knowledge}{%
\subsection{Pramanas (means of valid knowledge)}\label{p1-06--pramanas-means-of-valid-knowledge}}

\diagramhere{figures/p1-06-00}

\begingroup\small

\begin{longtable}[]{@{}
  >{\raggedright\arraybackslash}p{(\columnwidth - 4\tabcolsep) * \real{0.2723}}
  >{\raggedright\arraybackslash}p{(\columnwidth - 4\tabcolsep) * \real{0.1733}}
  >{\raggedright\arraybackslash}p{(\columnwidth - 4\tabcolsep) * \real{0.0624}}@{}}
\toprule\noalign{}
\begin{minipage}[b]{\linewidth}\raggedright
\textbf{School}
\end{minipage} & \begin{minipage}[b]{\linewidth}\raggedright
\textbf{Pramanas accepted}
\end{minipage} & \begin{minipage}[b]{\linewidth}\raggedright
\textbf{Count}
\end{minipage} \\
\midrule\noalign{}
\endhead
\bottomrule\noalign{}
\endlastfoot
Charvaka & Pratyaksha only & 1 \\
Buddhism, Vaisheshika & Pratyaksha, Anumana & 2 \\
Samkhya, Yoga, Jainism & + Shabda & 3 \\
\textbf{Nyaya} & + Upamana & 4 \\
Prabhakara Mimamsa & + Arthapatti & 5 \\
Bhatta Mimamsa, Advaita Vedanta & + Anupalabdhi & 6 \\
\end{longtable}

\endgroup

\begin{itemize}
\tightlist
\item
  \textbf{Arthapatti}: "Devadatta is fat but doesn\textquotesingle t eat by day ⇒ he eats at night."
\item
  \textbf{Anupalabdhi}: knowing the absence of a jar by not perceiving it.
\end{itemize}

\needspace{133pt}

\hypertarget{p1-06--anumana--five-membered-syllogism-panchavayava}{%
\subsection{Anumana --- Five-membered syllogism (Panchavayava)}\label{p1-06--anumana--five-membered-syllogism-panchavayava}}

\begin{asciiblock}{78pt}
\begin{Verbatim}[fontsize=\fontsize{8.8}{10.7}\selectfont]
1. Pratijna   (Proposition)  : The hill has fire.
2. Hetu       (Reason)       : Because it has smoke.
3. Udaharana  (Example)      : Whatever has smoke has fire, like a kitchen.
4. Upanaya    (Application)  : The hill has smoke (pervaded by fire).
5. Nigamana   (Conclusion)   : Therefore the hill has fire.
\end{Verbatim}
\end{asciiblock}

Terms: \textbf{Paksha} (minor term: hill), \textbf{Sadhya} (major term: fire), \textbf{Hetu/Linga} (middle term: smoke).
\textbf{Vyapti} = invariable concomitance (universal relation) between hetu and sadhya --- the basis of inference.

Kinds of anumana:

\begin{itemize}
\tightlist
\item
  By purpose: \textbf{Svarthanumana} (for oneself) and \textbf{Pararthanumana} (for others, uses 5 members).
\item
  By Nyaya (Gautama): \textbf{Purvavat} (cause → effect), \textbf{Sheshavat} (effect → cause), \textbf{Samanyatodrishta} (based on general observation).
\item
  By vyapti: Kevalanvayi, Kevalavyatireki, Anvaya-vyatireki.
\end{itemize}

\needspace{135pt}

\hypertarget{p1-06--hetvabhasa-fallacies-of-inference--5}{%
\subsection{Hetvabhasa (fallacies of inference --- 5)}\label{p1-06--hetvabhasa-fallacies-of-inference--5}}

\begingroup\small

\begin{longtable}[]{@{}
  >{\raggedright\arraybackslash}p{(\columnwidth - 2\tabcolsep) * \real{0.2189}}
  >{\raggedright\arraybackslash}p{(\columnwidth - 2\tabcolsep) * \real{0.5552}}@{}}
\toprule\noalign{}
\begin{minipage}[b]{\linewidth}\raggedright
\textbf{Hetvabhasa}
\end{minipage} & \begin{minipage}[b]{\linewidth}\raggedright
\textbf{Meaning}
\end{minipage} \\
\midrule\noalign{}
\endhead
\bottomrule\noalign{}
\endlastfoot
\textbf{Savyabhichara} (Anaikantika) & Irregular / inconclusive middle --- hetu found with and without sadhya \\
\textbf{Viruddha} & Contradictory --- hetu proves the opposite \\
\textbf{Satpratipaksha} & Counter-balanced --- another hetu proves the opposite \\
\textbf{Asiddha} (Sadhyasama) & Unproved --- hetu itself not established \\
\textbf{Badhita} (Kalatita) & Contradicted by other pramana (e.g. "fire is cold because it is a substance") \\
\end{longtable}

\endgroup

\needspace{194pt}

\hypertarget{p1-06--8-deeper-dive--venn-diagram-practice-conditionals--more-indian-logic}{%
\section{8. Deeper Dive --- Venn Diagram Practice, Conditionals \& More Indian Logic}\label{p1-06--8-deeper-dive--venn-diagram-practice-conditionals--more-indian-logic}}

\needspace{154pt}

\hypertarget{p1-06--81-six-standard-venn-situations}{%
\subsection{8.1 Six Standard Venn Situations}\label{p1-06--81-six-standard-venn-situations}}

\begin{asciiblock}{99pt}
\begin{Verbatim}[fontsize=\fontsize{8.8}{10.7}\selectfont]
 All A are B        No A is B          Some A are B
 ┌───────────┐      ┌─────┐ ┌─────┐    ┌──────┬──┬──────┐
 │ B  ┌───┐  │      │  A  │ │  B  │    │  A   │AB│   B  │
 │    │ A │  │      └─────┘ └─────┘    └──────┴──┴──────┘
 │    └───┘  │
 └───────────┘
 Some A are not B: the part of A outside B is non-empty (shade / mark with x)
\end{Verbatim}
\end{asciiblock}

\textbf{Method for "Statements → Conclusions"}:

\begin{enumerate}
\def\labelenumi{\arabic{enumi}.}
\tightlist
\item
  Draw the \textbf{minimum} case allowed by each statement.
\item
  A conclusion follows only if it holds in \textbf{every} possible diagram.
\item
  "Some A are not B" from "All A are B"? No. "Some B are A" from "All A are B"? Yes (conversion by limitation).
\end{enumerate}

\textbf{Worked}: Statements: All pens are books. Some books are bags.

\begin{itemize}
\tightlist
\item
  Conclusion I: Some pens are bags --- \textbf{does not follow} (bags may touch only the books outside pens).
\item
  Conclusion II: Some bags are books --- \textbf{follows} (conversion of "Some books are bags").
\item
  "Either I or II" type options apply only when two conclusions form a complementary pair (e.g. Some A are B / No A is B).
\end{itemize}

\needspace{135pt}

\hypertarget{p1-06--82-immediate-inference}{%
\subsection{8.2 Immediate Inference}\label{p1-06--82-immediate-inference}}

\begingroup\small

\begin{longtable}[]{@{}
  >{\raggedright\arraybackslash}p{(\columnwidth - 4\tabcolsep) * \real{0.1956}}
  >{\raggedright\arraybackslash}p{(\columnwidth - 4\tabcolsep) * \real{0.2689}}
  >{\raggedright\arraybackslash}p{(\columnwidth - 4\tabcolsep) * \real{0.3748}}@{}}
\toprule\noalign{}
\begin{minipage}[b]{\linewidth}\raggedright
\textbf{Operation}
\end{minipage} & \begin{minipage}[b]{\linewidth}\raggedright
\textbf{From}
\end{minipage} & \begin{minipage}[b]{\linewidth}\raggedright
\textbf{To}
\end{minipage} \\
\midrule\noalign{}
\endhead
\bottomrule\noalign{}
\endlastfoot
Conversion & E: No S is P / I: Some S are P & No P is S / Some P are S (valid) \\
Conversion by limitation & A: All S are P & Some P are S \\
Obversion & All S are P & No S is non-P (change quality, negate predicate) \\
Contraposition & All S are P & All non-P are non-S \\
\end{longtable}

\endgroup

O-propositions cannot be converted.

\needspace{133pt}

\hypertarget{p1-06--83-conditional-hypothetical-statements}{%
\subsection{8.3 Conditional (Hypothetical) Statements}\label{p1-06--83-conditional-hypothetical-statements}}

\begin{asciiblock}{78pt}
\begin{Verbatim}[fontsize=\fontsize{8.8}{10.7}\selectfont]
Valid:    If P then Q; P; ∴ Q          (Modus ponens)
Valid:    If P then Q; not Q; ∴ not P  (Modus tollens)
Invalid:  If P then Q; Q; ∴ P          (Affirming the consequent)
Invalid:  If P then Q; not P; ∴ not Q  (Denying the antecedent)
"P only if Q" ≡ If P then Q    "P unless Q" ≡ If not Q then P
\end{Verbatim}
\end{asciiblock}

\needspace{135pt}

\hypertarget{p1-06--84-truth--validity-combinations}{%
\subsection{8.4 Truth \& Validity Combinations}\label{p1-06--84-truth--validity-combinations}}

\begingroup\small

\begin{longtable}[]{@{}
  >{\raggedright\arraybackslash}p{(\columnwidth - 4\tabcolsep) * \real{0.0934}}
  >{\raggedright\arraybackslash}p{(\columnwidth - 4\tabcolsep) * \real{0.1043}}
  >{\raggedright\arraybackslash}p{(\columnwidth - 4\tabcolsep) * \real{0.2883}}@{}}
\toprule\noalign{}
\begin{minipage}[b]{\linewidth}\raggedright
\textbf{Premises}
\end{minipage} & \begin{minipage}[b]{\linewidth}\raggedright
\textbf{Conclusion}
\end{minipage} & \begin{minipage}[b]{\linewidth}\raggedright
\textbf{Can the argument be valid?}
\end{minipage} \\
\midrule\noalign{}
\endhead
\bottomrule\noalign{}
\endlastfoot
True & True & Yes \\
True & False & \textbf{No} --- impossible for a valid argument \\
False & True & Yes \\
False & False & Yes \\
\end{longtable}

\endgroup

\needspace{100pt}

\hypertarget{p1-06--85-more-indian-logic}{%
\subsection{8.5 More Indian Logic}\label{p1-06--85-more-indian-logic}}

\begin{itemize}
\tightlist
\item
  \textbf{Prama} = valid knowledge; \textbf{Aprama} = invalid knowledge (doubt --- samshaya, error --- viparyaya, hypothetical --- tarka).
\item
  \textbf{Pratyaksha} types (Nyaya): \textbf{Nirvikalpaka} (indeterminate) and \textbf{Savikalpaka} (determinate); also laukika (ordinary) and alaukika (extraordinary).
\item
  \textbf{Upamana}: knowledge of a word--object relation through similarity (learning that a gavaya resembles a cow).
\item
  \textbf{Shabda}: testimony of a reliable person (\textbf{apta vakya}); Vedic and secular.
\item
  \textbf{Jain logic}: \textbf{Anekantavada} (many-sidedness), \textbf{Syadvada} --- 7 modes of predication (\textbf{Saptabhangi naya}): syad asti, syad nasti, syad asti-nasti, syad avaktavya, \ldots{}
\item
  \textbf{Buddhist logic}: Dignaga and Dharmakirti; accept only perception and inference.
\item
  \textbf{Vyapti} is established through \textbf{anvaya} (agreement in presence: where smoke, there fire) and \textbf{vyatireka} (agreement in absence: where no fire, no smoke).
\item
  Nyaya Sutra author: \textbf{Gautama (Akshapada)}; 16 categories (padarthas) beginning with pramana and prameya.
\end{itemize}

\needspace{135pt}

\hypertarget{p1-06--previous-year-questions-pyq-pattern}{%
\section{Previous Year Questions (PYQ pattern)}\label{p1-06--previous-year-questions-pyq-pattern}}

\textbf{Square of opposition}

\begin{enumerate}
\def\labelenumi{\arabic{enumi}.}
\tightlist
\item
  If "All poets are dreamers" is true, which is false? → \textbf{"Some poets are not dreamers" (O)} and "No poets are dreamers" (E)
\item
  If "Some students are not intelligent" (O) is false, then "All students are intelligent" (A) is: \textbf{True}
\item
  Two propositions which cannot both be true but can both be false: \textbf{Contraries (A--E)}
\item
  Two propositions which cannot both be false but can both be true: \textbf{Sub-contraries (I--O)}
\item
  Which proposition distributes only the predicate? \textbf{O}
\end{enumerate}

\textbf{Syllogism}

\begin{enumerate}
\def\labelenumi{\arabic{enumi}.}
\setcounter{enumi}{5}
\tightlist
\item
  Premises: All cats are animals. Some animals are dogs. Conclusion "Some cats are dogs" → \textbf{Does not follow} (undistributed middle)
\item
  Statements: No A is B. All C are B. Conclusion: \textbf{No C is A} --- follows.
\item
  Mood and figure of: "All M are P; All S are M; ∴ All S are P" → \textbf{AAA-1 (Barbara)}
\end{enumerate}

\textbf{Deductive/inductive}

\begin{enumerate}
\def\labelenumi{\arabic{enumi}.}
\setcounter{enumi}{8}
\tightlist
\item
  In a deductive argument, the conclusion: \textbf{follows necessarily}
\item
  An argument which is valid and has true premises is: \textbf{Sound}
\item
  Inductive arguments are evaluated as: \textbf{strong/weak}
\end{enumerate}

\textbf{Fallacies}

\begin{enumerate}
\def\labelenumi{\arabic{enumi}.}
\setcounter{enumi}{11}
\tightlist
\item
  "He is a criminal, so his argument about taxes is wrong" --- \textbf{Ad hominem}
\item
  "No one has proved ghosts don\textquotesingle t exist, so they exist" --- \textbf{Ad ignorantiam}
\item
  "I wore a lucky shirt and won --- so the shirt causes wins" --- \textbf{Post hoc (false cause)}
\end{enumerate}

\textbf{Language}

\begin{enumerate}
\def\labelenumi{\arabic{enumi}.}
\setcounter{enumi}{14}
\tightlist
\item
  Denotation of a term refers to: \textbf{the objects to which it applies}
\item
  "Close the door" serves which function? \textbf{Directive}
\end{enumerate}

\textbf{Indian logic}

\begin{enumerate}
\def\labelenumi{\arabic{enumi}.}
\setcounter{enumi}{16}
\tightlist
\item
  Number of pramanas accepted by Nyaya: \textbf{4}
\item
  Charvakas accept only: \textbf{Perception}
\item
  Advaita Vedanta accepts how many pramanas? \textbf{6}
\item
  In "the hill has fire because it has smoke", the middle term (hetu) is: \textbf{smoke}; paksha is: \textbf{hill}
\item
  The invariable relation between hetu and sadhya is: \textbf{Vyapti}
\item
  Inference for others is: \textbf{Pararthanumana}
\item
  Arrange Panchavayava: \textbf{Pratijna, Hetu, Udaharana, Upanaya, Nigamana}
\item
  A hetu that proves the contrary of what is intended: \textbf{Viruddha}
\item
  Knowledge of non-existence is obtained through: \textbf{Anupalabdhi}
\item
  Arthapatti is accepted by: \textbf{Mimamsa (Prabhakara \& Bhatta) and Advaita}
\end{enumerate}

\textbf{Analogy}

\begin{enumerate}
\def\labelenumi{\arabic{enumi}.}
\setcounter{enumi}{26}
\tightlist
\item
  Book : Author :: Statue : \textbf{Sculptor}
\end{enumerate}

\textbf{More practice questions}

\begin{enumerate}
\def\labelenumi{\arabic{enumi}.}
\setcounter{enumi}{27}
\tightlist
\item
  Obverse of "All S are P": \textbf{No S is non-P}
\item
  Which proposition cannot be converted? \textbf{O}
\item
  Converse of "All doctors are graduates" by limitation: \textbf{Some graduates are doctors}
\item
  A valid argument cannot have: \textbf{true premises and a false conclusion}
\item
  "If it rains, the match is cancelled. The match is cancelled. So it rained." --- \textbf{affirming the consequent}
\item
  Statements: No cat is a dog. All dogs are animals. Conclusion "Some animals are not cats" --- \textbf{follows}
\item
  Syadvada is associated with: \textbf{Jainism}
\item
  Number of modes in Saptabhangi: \textbf{7}
\item
  Determinate perception in Nyaya: \textbf{Savikalpaka}
\item
  Vyapti by agreement in absence: \textbf{Vyatireka}
\item
  Author of Nyaya Sutras: \textbf{Gautama (Akshapada)}
\item
  Knowledge from similarity (gavaya--cow): \textbf{Upamana}
\item
  Analogy: Doctor : Hospital :: Teacher : \textbf{School}
\end{enumerate}

\begin{revbox}

\begin{itemize}
\tightlist
\item
  A(S) E(S,P) I(none) O(P) --- distribution
\item
  Contradictory A--O, E--I · Contrary A--E · Sub-contrary I--O
\item
  Charvaka 1 · Buddhist 2 · Samkhya 3 · Nyaya 4 · Prabhakara 5 · Bhatta/Advaita 6
\item
  5 Hetvabhasa: Savyabhichara, Viruddha, Satpratipaksha, Asiddha, Badhita
\end{itemize}

\end{revbox}
